\section{Marginal Contribution of Feature Buckets and Individual Features}\label{sec:marginal}

Section~\ref{sec:descriptive} established the panel, the target, and the evaluation
protocol; it also fixed the reference point against which everything in this paper is
measured, the per-bar least-squares fit of the HAR ladder plus calendar dummies. This
section asks the first substantive question: \emph{where does forecast skill actually come
from?} We answer it at two levels of resolution --- the exogenous feature buckets defined
in Section~\ref{sec:descriptive}, and, one level down, the individual variables inside them
--- and we answer it deliberately within a single model class, ordinary least squares.

The restriction to least squares is the point of the exercise, not a limitation of it.
Attribution statements that are computed under a flexible learner confound two distinct
questions: whether a feature carries information, and whether a particular estimator
happens to be able to exploit it. By holding the estimator fixed at the simplest possible
linear map we obtain marginal contributions that are properties of the \emph{data}. What a
richer hypothesis class does with the same features is the subject of
Section~\ref{sec:linvsnonlin}, and the comparison between the two sections is only
interpretable because this one refuses to move the estimator.

\paragraph{What ``least squares'' means here, precisely.} Honesty requires that we state
the estimator used in each experiment rather than assert a uniform ``OLS'' that was never
run in that exact form. Three closely related configurations appear below.

\begin{enumerate}
\item \textbf{The battery's linear arms} (Section~\ref{sec:marginal} 3.1). The primary
attribution table is the linear slice of the same causal-tuning battery that
Section~\ref{sec:linvsnonlin} analyses in full: rolling least squares over a window of
$24{,}000$ bars ($500$ trading days), refit at every bar by exact rank-one
downdate/update of the normal equations (reproducing the batch solution to
$1.8 \times 10^{-11}$), with a ridge \citep{HoerlKennard1970}, lasso
\citep{Tibshirani1996} or elastic-net \citep{ZouHastie2005} penalty whose strength is
re-selected causally from a forward validation split. The penalties selected are small,
and the ridge column is the working least-squares surrogate throughout. The battery's own
\texttt{baseline} row bounds the distance from exact OLS directly: penalizing the HAR
ladder with no exogenous features present moves QLIKE by only $+0.0003$ against the
$\alpha = 0$ incumbent --- and \emph{in the wrong direction}, which is itself evidence
that shrinkage is doing nothing for these tables beyond numerical stabilisation.
\item \textbf{HAR-unpenalized, exog-penalized every-bar fits}
(Sections~\ref{sec:marginal} 3.2 and~\ref{sec:marginal} 3.5). The six HAR columns enter
\emph{without} any penalty, which by the Frisch--Waugh--Lovell theorem
\citep{FrischWaugh1933,Lovell1963} makes the HAR block an exact OLS fit on the residualised
problem; only the exogenous block carries a penalty. The baseline row of every such design
is therefore literally OLS.
\item \textbf{Ridge at $\alpha = 1$ on the raw research panel}
(Section~\ref{sec:marginal} 3.4), used for the exhaustive $2^8$ factorial, where $256$
cells $\times$ two training windows had to be run every-bar.
\end{enumerate}

\noindent A literal $\alpha = 0$ battery with exogenous features has \emph{not} been run,
and we do not claim one. The incumbent code path accepts $\alpha = 0$ directly, so this is
a pending robustness run rather than an open question of feasibility; we flag it again in
the limitations at the end of the section. Nothing below depends on the distinction at the
displayed precision, but the reader is entitled to know which fit produced which number.

\paragraph{A standing warning about levels.} Three mutually incompatible QLIKE level
regimes appear in this paper, because the underlying panels differ in construction and in
evaluation window. The diurnally adjusted panel used for the headline results sits at
$\approx 0.126$--$0.134$; the COVID-inclusive every-bar battery sits at
$\approx 0.142$--$0.149$; the raw research panel used for the factorial, which applies no
diurnal adjustment and retains overnight bars, sits at $\approx 0.35$--$0.79$. \emph{Levels
are never comparable across these three regimes.} Every table below carries its panel and
its evaluation window in the caption, and every claim we draw is a within-table ranking, a
within-table difference, or a test statistic.

\subsection{Bucket-Level Marginal Contribution}

The primary attribution experiment adds each exogenous bucket, one at a time, to the HAR
ladder, and scores the resulting forecast on the full battery panel of
Section~\ref{sec:descriptive}. Table~\ref{tab:marginal_buckets_wf} is the linear slice of
the battery: the same rows, the same panel, and the same incumbent as every other
bucket-by-model table in this paper --- only the model columns differ. Here they are the
three members of the least-squares class; Section~\ref{sec:linvsnonlin} will widen the
same table to tree ensembles, and Section~\ref{sec:algorithms} to nonparametric
estimators, without moving anything else.

% Provenance: writeup/metrics_table_causal_tune_plus_spectral.csv (family=causal-tuned,
% estimators ridge/reclasticnet/reclasso), identical source as Section 4's battery table;
% incumbent from writeup/incumbent_metrics.json.
\begin{table}[H]
\centering
\small
\setlength{\tabcolsep}{4pt}
\caption{Marginal contribution of each exogenous bucket: the linear slice of the battery.
\textbf{Panel: causal-tuning battery,} diurnally adjusted, $n = 218{,}934$ thirty-minute
bars, identical rows in every cell; incumbent QLIKE $= 0.13415$ (per-bar OLS on HAR $+$
calendar). \textbf{Estimators:} rolling least squares, window $24{,}000$ bars ($500$
trading days), refit every bar by exact rank-one updates, penalty re-selected causally.
Cells report QLIKE with the Diebold--Mariano $t$-statistic against the incumbent in
parentheses; negative $t$ favors the arm. Best cell per row in bold. The
\texttt{all\_features} row stacks all $41$ exogenous series
($\approx 526$ columns after the moving-average ladder and availability indicators).}
\label{tab:marginal_buckets_wf}
\begin{tabular}{@{}lccc@{}}
\toprule
Bucket & Ridge & Elastic net & Lasso \\
\midrule
\texttt{all\_features} & \textbf{0.12788} ($-19.9$) & 0.12827 ($-19.3$) & 0.12907 ($-13.2$) \\
\texttt{moments}       & \textbf{0.13085} ($-19.2$) & 0.13114 ($-17.7$) & 0.13179 ($-8.9$) \\
\texttt{liquidity}     & \textbf{0.13120} ($-13.4$) & 0.13184 ($-11.5$) & 0.13206 ($-11.2$) \\
\texttt{implied\_vol}  & \textbf{0.13252} ($-11.3$) & 0.13262 ($-10.6$) & 0.13274 ($-9.8$) \\
\texttt{market\_vw}    & \textbf{0.13291} ($-9.1$)  & 0.13316 ($-7.3$)  & 0.13321 ($-6.4$) \\
\texttt{market\_ew}    & \textbf{0.13319} ($-7.2$)  & 0.13344 ($-5.4$)  & 0.13353 ($-4.6$) \\
\texttt{vol\_demand}   & \textbf{0.13337} ($-4.1$)  & 0.13361 ($-3.4$)  & 0.13359 ($-3.3$) \\
\texttt{sentiment}     & \textbf{0.13347} ($-5.7$)  & 0.13377 ($-3.5$)  & 0.13390 ($-2.1$) \\
\midrule
\texttt{baseline} (no exog) & \textbf{0.13445} ($+5.0$) & 0.13471 ($+7.5$) & 0.13475 ($+7.0$) \\
Incumbent: OLS on \texttt{baseline} & \multicolumn{3}{c}{0.13415} \\
\bottomrule
\end{tabular}
\end{table}

Every comparison in Table~\ref{tab:marginal_buckets_wf} is bar-for-bar against the same
incumbent on the same rows, so each cell carries its own significance statement. Three
facts follow, and they organise the rest of the paper.

\paragraph{Every bucket improves on HAR, significantly.} Reading down the ridge column,
there is no exogenous bucket whose addition degrades the forecast, and no bucket whose
improvement is attributable to sampling noise: the weakest arm in the table,
\texttt{vol\_demand}, still rejects equal predictive accuracy at $t = -4.1$. The spread
between the best single bucket (\texttt{moments}, $-0.00331$ against the incumbent) and
the weakest (\texttt{sentiment}, $-0.00068$) is a factor of five, but the sign is
uniform. This uniformity is itself a finding, and a hard-won one: in an earlier revision
of the pipeline, three of these buckets --- \texttt{sentiment}, \texttt{implied\_vol} and
\texttt{vol\_demand} --- appeared to \emph{lose} to HAR, and were provisionally written up
as such. Each of those apparent losses turned out to be a scaling defect in the feature
construction rather than an absence of signal, and each bucket became a genuine
improvement once the defect was repaired. The general lesson, which we state because it
bears on how any such table should be read, is that a feature that appears to hurt is
very often a feature the pipeline mis-scales. An earlier, independent walk-forward
battery at a fixed $\alpha = 1$ penalty --- a different configuration on the same panel
--- produces the same ranking with \texttt{moments} first among single buckets and the
joint set dominant, so the attribution is not an artefact of the causal penalty
selection.
% Replication source: writeup/session_notes_2026-06-23_feature_investigation.md,
% "Final validated ranking (real walk-forward QLIKE, full 218,934, vs HAR 0.13460)".

\paragraph{The buckets overlap heavily: contribution is strongly sublinear.} Summing the
seven individual bucket improvements in the ridge column gives
\[
\textstyle\sum_b \Delta_b
= -(0.00331 + 0.00295 + 0.00163 + 0.00125 + 0.00096 + 0.00078 + 0.00068)
= -0.01155 ,
\]
whereas the joint \texttt{all\_features} cell delivers only $-0.00627$ --- barely more
than half of the additive prediction ($-0.01155 / -0.00627 = 1.84$). The buckets are
therefore not seven independent channels of information about future variance; they are
seven partially redundant views of a much smaller amount of information. Note that the
joint cell nonetheless beats the best single bucket by a further $-0.00297$, so the
redundancy is partial rather than total: stacking still pays, but it pays at roughly half
the rate a naive addition of marginal gains would suggest. (An independent, test-set-tuned
oracle table computed at a different training window and on a different evaluation slice
--- diagnostic only, and not level-comparable to Table~\ref{tab:marginal_buckets_wf} ---
reproduces the same qualitative sublinearity, with a joint improvement smaller than the sum
of its singles.) Sublinearity of this magnitude is the single most important structural
fact about the feature set, and it recurs at the individual-feature level in
Section~\ref{sec:marginal} 3.4 below.

\paragraph{A provenance caution.} The \texttt{market\_ew} bucket is worth one sentence of
warning about how easily an attribution table can be corrupted upstream of the estimator.
In the fixed-$\alpha$ replication battery just cited, the same bucket evaluated on the
same panel with the same estimator but with an overnight fill routine that imputed a
constant $1.0$ into the cross-sectional return moments scores $\Delta = -0.00044$ against
its HAR reference rather than the clean $-0.00128$: the imputation artefact silently
destroyed roughly two thirds of the bucket's genuine contribution. Marginal-contribution
tables are only as trustworthy as the missing-data handling beneath them, and the reader
should treat the imputation policy as part of the estimator specification rather than as
a pre-processing detail.

\subsection{Robustness Across Evaluation Windows, and the Model Confidence Set}

The Diebold--Mariano statistics of Table~\ref{tab:marginal_buckets_wf}
\citep{DieboldMariano1995} are computed as the $t$-ratio of the mean per-bar loss
differential to its Newey--West HAC standard error with the Harvey--Leybourne--Newbold
small-sample correction \citep{HarveyLeybourneNewbold1997}. Two questions remain that the
battery alone cannot answer: whether the bucket attributions survive a different
evaluation window and protocol, and whether any \emph{single} bucket is statistically
defensible against the joint set. Both are answered by an independent every-bar
robustness design, which we summarise in prose because its QLIKE levels live on a
different panel and are not comparable to the battery: a rolling
$1{,}000$-trading-day window, HAR block unpenalized (exact OLS via Frisch--Waugh--Lovell)
with only the exogenous block penalized, scored on a fixed COVID-inclusive window from
January 2018 to May 2025, $n = 74{,}934$ bars over $1{,}843$ trading days with identical
rows in every cell.
% Provenance: writeup/beat4_table_with_trees.tex (linear rows) and
% writeup/penalty_estimator_floor_2026-07-01.md section 7 (DM / MCS).

The robustness design reproduces the battery's qualitative structure in full. Every
bucket again rejects equal predictive accuracy against HAR-only at any conventional
level --- the weakest, \texttt{vol\_demand}, at $t = -6.2$
($p = 5.0 \times 10^{-10}$), the joint set at $t = -13.0$ --- and the joint set again
dominates every single bucket. With $74{,}934$ paired loss observations this is
unsurprising and should not be over-read: the tests establish that the improvements are
not sampling noise, not that they are economically large. The bucket \emph{ranking}
rotates with the window --- \texttt{implied\_vol} leads there, \texttt{moments} in
Table~\ref{tab:marginal_buckets_wf} --- which is exactly what one expects when the
evaluation window changes to one dominated by a volatility crisis: implied-volatility
levels and liquidity carry the most information through COVID. The rotation is a further
argument against promoting any single bucket.

The model confidence set delivers the sharper statement. At the $90\%$ level, and at
\emph{every} training window tested from $1{,}000$ to $3{,}500$ trading days, the
bucket-level MCS is the singleton $\{\texttt{all\_buckets}\}$: every individual bucket is
eliminated. The joint set against HAR-only yields Diebold--Mariano statistics ranging from
$-12.8$ to $-16.2$ across those windows, with a mean loss differential of $-0.0054$.

The conclusion is worth stating without hedging. Individual buckets are each significantly
better than HAR, and they are each significantly \emph{worse} than the union of all
buckets, uniformly and at every window we examined. There is no defensible ``top bucket''
in a statistical sense --- only a defensible joint feature set. Any narrative that promotes
one bucket as the source of skill is reading a point estimate that the significance
machinery does not support. This is the natural companion to the sublinearity finding of
the previous subsection: the buckets overlap so heavily that no proper subset survives, yet
each contributes enough unique information that dropping it is detectable.

\subsection{A Block-Level Type-III Decomposition}

The bucket batteries treat the exogenous features as one undifferentiated addition to HAR.
A finer question is how the deployed feature matrix partitions into functionally distinct
\emph{blocks}, and what each block is worth net of all the others.

The natural instrument is the Frisch--Waugh--Lovell theorem
\citep{FrischWaugh1933,Lovell1963}. In a least-squares fit, the coefficient on a block of
regressors is identical to the coefficient obtained by first residualising both that block
and the dependent variable on all remaining regressors, and then regressing one residual on
the other. The block's contribution is therefore exactly its contribution \emph{orthogonal
to everything else in the model} --- a Type-III, leave-one-block-out decomposition rather
than a sequential Type-I one that would depend on the order in which blocks are entered.
Operationally we implement this by nested causal walk-forward refits: each block's unique
contribution is the degradation in out-of-sample QLIKE when that block alone is removed
from the full specification, with all remaining blocks refit optimally. Because the blocks
are strongly collinear, refitting matters; decomposing a single fitted coefficient vector
would attribute shared variance arbitrarily.

% Provenance: writeup/mechanism_and_data_to_buy_2026-06-28.md, section 3
% "FWL block attribution of the linear base (fwl_attribution.py; full fit reproduces 0.12314)".
\begin{table}[H]
\centering
\caption{Frisch--Waugh--Lovell block decomposition of the deployed linear base.
\textbf{Panel:} diurnally adjusted full out-of-sample panel; the full fit reproduces QLIKE
$0.12314$. Each entry is the block's unique (Type-III) contribution, i.e.\ the
out-of-sample QLIKE degradation from removing that block with all others refit. Computed
under the production fixed-$\alpha$ elastic net on the exogenous block, so the
orthogonalisation is approximate rather than exact --- see the caveat in the text.}
\label{tab:marginal_fwl}
\begin{tabular}{llr}
\toprule
block & content & unique $\Delta$QLIKE \\
\midrule
HAR    & the six-lag realised-variance ladder                     & $-0.02208$ \\
EXOG   & 359 raw exogenous columns                                & $-0.00449$ \\
REGIME & HAR $\times$ \{open, close\} session-edge interactions   & $-0.00225$ \\
CUMRV  & \texttt{cumrv\_x\_close}, one engineered path feature    & $-0.00122$ \\
\bottomrule
\end{tabular}
\end{table}

The HAR ladder dominates: it accounts for $96\%$ of the explainable variance of the full
model, in the sense that HAR alone reaches $R^2 = 0.594$ against the full specification's
$0.617$. This is the quantitative form of the observation, familiar since
\citet{Corsi2009}, that multi-scale volatility persistence is most of what there is to
know about future realised variance at these horizons. Everything built on top of it ---
the entire exogenous data programme, the session-edge interactions, the engineered path
feature --- shares a residual roughly one twentieth the size of the HAR contribution.

The more interesting comparison is per feature. The CUMRV block is a \emph{single} column,
worth $-0.00122$ on its own; the EXOG block is $359$ columns worth $-0.00449$ jointly, or
about $1.3 \times 10^{-5}$ per column. One mechanism-targeted engineered feature is thus
worth nearly two orders of magnitude more, per feature, than the average raw exogenous
series. We regard this as the single sharpest piece of evidence in the paper about where
effort should go: a feature that directly measures the mechanism (here, the accumulation of
volatility along the intraday path, gated to the close and after-hours session) dominates
the strategy of adding more loosely related series and letting the penalty sort them out.
The corollary, developed in Section~\ref{sec:linvsnonlin}, is that the binding constraint
is information rather than estimator capacity.

\paragraph{Caveat.} These numbers were computed under the production fixed-$\alpha$ elastic
net on the exogenous block, not under exact least squares. Frisch--Waugh--Lovell is exact
in OLS; under a penalty the residualisation is only approximate, because the penalty is not
invariant to the reparameterisation that the theorem performs, and shared variance can be
reallocated between blocks by the shrinkage rather than by the data. The direction and
rough magnitude of the decomposition are robust --- the HAR block's dominance is not a
$0.001$-scale effect --- but the exact figures should be re-derived under $\alpha = 0$, and
we flag that rerun alongside the pending unpenalized bucket battery in the limitations
below.

\subsection{Individual Features: An Exact $2^8$ Factorial}

Bucket-level attribution cannot say which \emph{variables} matter. To answer that we ran a
complete $2^8$ factorial over the eight variable bundles that survived a first-stage
screen, where a bundle is one variable's six-lag moving-average ladder together with its
availability indicator. All $256$ cells were evaluated, so the main effects and
interactions reported below are exact averages over the design rather than estimates from a
fitted response surface.

This experiment lives on the \emph{raw research panel}: no diurnal adjustment, no rank
transform, overnight bars included. Its QLIKE levels are consequently in the
$0.35$--$0.79$ range and \textbf{cannot be compared to any other table in this paper}. The
bundle gains here include intraday seasonality that the adjusted panel removes by
construction. Rankings, main effects, and Diebold--Mariano statistics are the deliverable.
Two training windows were run, a primary $tw = 48{,}000$ bars ($1{,}000$ trading days;
evaluation 9 October 2008 -- 30 April 2024, $n = 198{,}059$) and a replicate
$tw = 144{,}000$ ($3{,}000$ days; evaluation 25 April 2016 -- 30 April 2024,
$n = 102{,}059$); the two windows have different evaluation spans and their levels are not
comparable to each other either, so the replicate serves as a stability check on signs and
ranks.

% Provenance: writeup/beat4_meeting_2026-07-02.tex, Table 2 (single-bundle cells,
% exact factorial main effects and two-way interactions in the caption).
\begin{table}[H]
\centering
\caption{Individual-feature factorial. \textbf{Panel:} raw research build (no diurnal
adjustment or rank transform, overnight bars retained); ridge $\alpha = 1$ on the
square-root target with the RV-HAR ladder and intraday hour dummies always included;
refit every bar. \textbf{Levels are not comparable to any other table in this paper, nor
between the two windows.} ME is the exact factorial main effect on QLIKE, mean(bundle on)
$-$ mean(bundle off) over all $256$ cells. Single-cell QLIKE is the bundle added alone to
the baseline; DM is the Diebold--Mariano $t$-statistic against the baseline row of the same
window.}
\label{tab:marginal_factorial}
\begin{tabular}{lrrrrrr}
\toprule
& \multicolumn{3}{c}{$tw = 48{,}000$ (2008--2024)} & \multicolumn{3}{c}{$tw = 144{,}000$ (2016--2024)} \\
\cmidrule(lr){2-4}\cmidrule(l){5-7}
bundle & ME & QLIKE & DM & ME & QLIKE & DM \\
\midrule
\texttt{absret}    & $-0.199$ & 0.38950 & $-37.3$ & $-0.257$ & 0.38686 & $-58.8$ \\
\texttt{volume}    & $-0.108$ & 0.50567 & $-31.3$ & $-0.063$ & 0.57475 & $-44.9$ \\
\texttt{vvix}      & $-0.023$ & 0.68515 & $-12.2$ & $-0.014$ & 0.66325 & $-23.9$ \\
\texttt{sentcount} & $-0.011$ & 0.71939 & $-11.9$ & $+0.003$ & 0.69639 & $-7.1$ \\
\texttt{vdem\_all} & $-0.010$ & 0.72088 & $-14.7$ & $+0.000$ & 0.68742 & $-8.4$ \\
\texttt{attention} & $-0.010$ & 0.72629 & $-10.0$ & $+0.002$ & 0.68812 & $-8.4$ \\
\texttt{vdem\_spx} & $-0.009$ & 0.73552 & $-12.6$ & $+0.001$ & 0.67324 & $-13.0$ \\
\texttt{bipow}     & $-0.000$ & 0.78314 & $-6.2$  & $-0.000$ & 0.72567 & $-2.1$ \\
\midrule
baseline (HAR ladder $+$ hour dummies) & --- & 0.78614 & --- & --- & 0.72675 & --- \\
\bottomrule
\end{tabular}
\end{table}

\paragraph{One dominant feature, one partially redundant second, and a long thin tail.}
Absolute returns carry a main effect of $-0.199$ and $-0.257$ in the two windows, roughly
twice the next largest and an order of magnitude larger than everything below the top two.
Volume follows at $-0.108$ and $-0.063$. Below those two the main effects collapse to the
third decimal, and in the replicate window four of them ($\texttt{sentcount}$,
$\texttt{vdem\_all}$, $\texttt{attention}$, $\texttt{vdem\_spx}$) change sign to weakly
positive, which is to say they are indistinguishable from zero once the training window is
tripled. Bipower variation is a genuine null in both windows, $-0.000$ and $-0.000$: it is
the one bundle in the design that the factorial cannot distinguish from adding nothing at
all, and we report it as such rather than folding it into the tail. That an established
jump-robust volatility estimator adds no exact main effect here is informative --- against
a six-lag HAR ladder already fit on the same realised-variance series, bipower variation
appears to be spanned.

\paragraph{The two leading features are partly redundant with one another.} The largest
two-way interactions are $\texttt{absret} \times \texttt{volume}$ at $+0.094$ and $+0.064$,
and $\texttt{absret} \times \texttt{vvix}$ at $+0.020$ and $+0.010$. Under this sign
convention a positive interaction means the joint gain is \emph{less} than additive.
The $\texttt{absret} \times \texttt{volume}$ interaction recovers roughly $87\%$ of
volume's own main effect in the primary window --- adding volume once absolute returns are
present buys a small fraction of what it buys on its own. This is the individual-feature
image of the bucket-level sublinearity of Section~\ref{sec:marginal} 3.1, one level down and
measured exactly rather than by summation.

\paragraph{Sparse beats dense.} Because the design is complete we can identify the best
subset directly rather than by search. The four-bundle set
$\{\texttt{absret}, \texttt{volume}, \texttt{vdem\_all}, \texttt{bipow}\}$ attains QLIKE
$0.34598$ with $\mathrm{DM} = -40.9$ against the baseline, beating the all-eight
specification at $0.35156$ with $\mathrm{DM} = -39.4$; the ordering replicates in the
second window ($0.36599$ versus $0.37023$). A six-bundle set chosen on a different
criterion lands at $0.35237$, between the two. Adding the four weakest bundles to the
four-bundle winner therefore \emph{costs} accuracy: their marginal information is smaller
than the estimation variance they introduce, even under a penalised fit at $\alpha = 1$.

Taken together, the picture at the individual-feature level is the same one the bucket
tables paint: one cheap, dominant feature; a second that is largely redundant with it; and
a long tail of effects small enough that including them is a variance decision rather than
a bias decision. The information in this feature set is dense, weak, and heavily
overlapping --- the framing we carry into Section~\ref{sec:linvsnonlin}.

\subsection{What Least Squares Keeps: The Support Structure}

A final diagnostic asks which columns a sparse least-squares fit actually retains. We fit
the exogenous block with an $\ell_1$ penalty \citep{Tibshirani1996} at its oracle strength,
HAR held unpenalized, across ten sequential walk-forward blocks of the $523$-column
exogenous matrix, and record the support in each block.

% Provenance: results/lasso_support.log (driver lasso_support.py), summarised in
% writeup/penalty_estimator_floor_2026-07-01.md section 5.
\begin{table}[H]
\centering
\caption{Support of the oracle $\ell_1$ fit on the exogenous block, across ten sequential
walk-forward blocks. \textbf{Panel:} $523$ penalized exogenous columns, training window
$48{,}000$ bars. ``Ever kept'' counts columns with a non-zero coefficient in at least one
block; ``always kept'' counts columns non-zero in all ten; ``killed'' counts columns zero
in every block.}
\label{tab:marginal_support}
\begin{tabular}{lrrrr}
\toprule
family & columns & ever kept & always kept & killed \\
\midrule
availability indicators & 246 & 11 & 0 & 235 \\
\texttt{liquidity}      & 72  & 58 & 10 & 14 \\
\texttt{vol\_demand}    & 48  & 26 & 0 & 22 \\
\texttt{moments}        & 42  & 32 & 8 & 10 \\
\texttt{market\_ew}     & 36  & 26 & 1 & 10 \\
\texttt{market\_vw}     & 36  & 27 & 1 & 9 \\
\texttt{sentiment}      & 18  & 16 & 0 & 2 \\
\texttt{implied\_vol}   & 18  & 16 & 2 & 2 \\
calendar                & 6   & 6  & 6 & 0 \\
other                   & 1   & 1  & 1 & 0 \\
\midrule
total                   & 523 & 219 & 29 & 304 \\
\bottomrule
\end{tabular}
\end{table}

Of the $523$ exogenous columns, $304$ are zeroed in every block. Of those $304$ kills, $235$
--- some $77\%$ --- are availability indicators: only $11$ of the $246$ missing-data flags
ever carry signal. This has a direct estimator-selection consequence that we note here
because it recurs in Section~\ref{sec:algorithms}: a pure $\ell_2$ penalty
\citep{HoerlKennard1970} cannot set these columns to zero, only shrink them, so they inject
variance into a dense fit; $\ell_1$ removes them cleanly. Where a bucket contains little
such junk --- \texttt{implied\_vol}, whose $18$ columns are dense and almost all
informative --- the ordering reverses and $\ell_2$ is preferable. Penalty choice tracks the
junk fraction of the design matrix, not any deep property of the signal.

The temporal structure of the support is the more striking result. Of the $219$ columns ever
kept, only $29$ are kept in every block; the remaining $190$ --- $87\%$ of the ever-selected
set --- enter and leave as the window rolls forward. The stable core is small and
recognisable: the six calendar and clock features are retained in every single block, joined
by ten liquidity, eight moments and two implied-volatility columns. Everything else
flickers.

Two readings of this are consistent with the rest of the section. Under the first, the
flicker is estimation noise, in which case the marginal value of any individual exogenous
column is essentially unmeasurable at this sample size and the joint bucket is the only
meaningful unit of attribution --- exactly what the model confidence set concluded in
Section~\ref{sec:marginal} 3.2. Under the second, the flicker is genuine regime dependence,
in which case a single fixed linear map is the wrong object and the appropriate model is one
that varies its effective support over time. Either way, the operative description of this
feature set is a small unstable core plus a wide, shallow tail, which is the dense-but-weak
regime that Section~\ref{sec:linvsnonlin} takes as its starting point.

\subsection{Limitations}

Four qualifications apply to everything above, and we state them plainly.

\begin{enumerate}
\item \textbf{The exact $\alpha = 0$ battery is pending.} No literal unpenalized
least-squares run with exogenous features has been executed. Section~\ref{sec:marginal} 3.1
uses a negligible ridge penalty at a $24{,}000$-row window; Sections~\ref{sec:marginal} 3.2
and~\ref{sec:marginal} 3.5 hold the HAR block exactly unpenalized but penalise the
exogenous block. The incumbent code path accepts $\alpha = 0$ directly, so this is a
scheduled robustness run, not an unresolved question.
\item \textbf{The Frisch--Waugh--Lovell decomposition is penalty-approximate.}
Table~\ref{tab:marginal_fwl} was computed under the production fixed-$\alpha$ elastic net.
FWL is exact only in OLS; under shrinkage the orthogonalisation reallocates shared variance
in a way that is not purely data-determined. The decomposition should be re-derived at
$\alpha = 0$ alongside the battery above.
\item \textbf{No per-coefficient standard errors are reported anywhere.} All inference in
this section is on \emph{forecast losses} --- Diebold--Mariano tests and model confidence
sets on per-bar QLIKE --- not on coefficients. We make no claims about the significance of
any individual estimated coefficient, and none should be inferred from the support table.
\item \textbf{The factorial lives on a different panel in different units.}
Table~\ref{tab:marginal_factorial} is built on the raw research panel with no diurnal
adjustment and overnight bars retained. Its levels are not comparable to
Table~\ref{tab:marginal_buckets_wf} or
Table~\ref{tab:marginal_fwl}, and its main effects are in that panel's units. Only rankings,
signs, interaction structure and test statistics transfer.
\end{enumerate}
